\documentclass{article}
\usepackage{geometry}
\geometry{a4paper, margin=1in}
\usepackage[utf8]{inputenc}
\usepackage{amsmath}
\usepackage{graphicx}
\usepackage[numbers]{natbib}
\usepackage{hyperref}

\title{Complex QA and Language Models Hybrid Architectures Survey}
\author{Research Assistant in AI}

\begin{document}

\maketitle

\begin{abstract}
This paper reviews the state-of-the-art in language model architectures and strategies for complex question-answering (QA) with a focus on hybridization. Large Language Models (LLMs) leverage extensive public data for standard problems, but specific complex questions require specialized architectures, knowledge, and methods. Recent projects like ChatGPT and GALACTICA illustrate both the potential and limitations of LLMs in complex QA. We review required skills and evaluation techniques, and integrate findings from projects like BIG, BLOOM, and HELM. We discuss challenges such as domain adaptation, multi-step QA, long-form QA, data protection, multimodal search, hallucinations, explainability, and temporal reasoning. We analyze current solutions and promising research trends including hybrid LLM architectures, training strategies, human reinforcement learning, neuro-symbolic grounding, and program synthesis.
\end{abstract}

\tableofcontents

\section{Introduction}
Large Language Models (LLMs) have revolutionized the field of question-answering (QA) by demonstrating the ability to process and generate human-like text based on extensive training data. However, when faced with complex questions that go beyond the scope of their training data, their limitations become evident. Complex Question Answering (CQA) requires models not only to understand and generate natural language but also to apply domain-specific knowledge, reasoning, and often multi-step interactions to derive accurate answers.

This paper aims to survey the current state-of-the-art in hybrid architectures that combine the strengths of LLMs with specialized components to effectively tackle complex QA tasks. We will explore necessary skills and methodologies for evaluating these systems, delve into significant challenges, and analyze innovative solutions and trends in hybrid architectures.

\section{Skills and Evaluation Techniques}
To address complex QA, models need a range of skills including understanding context, reasoning, and applying domain-specific knowledge. Evaluation techniques for these systems need to focus on several aspects:

\subsection{Necessary Skills}
\begin{itemize}
    \item \textbf{Contextual Understanding}: The ability to understand the context of a question and its components is crucial for accurate answers.
    \item \textbf{Knowledge Integration}: Leveraging external knowledge sources, such as structured databases or domain-specific ontologies, enhances the model’s capacity to provide informed responses.
    \item \textbf{Reasoning and Inference}: Effective reasoning is required to draw conclusions from given data, especially in multi-step or multi-hop QA tasks.
    \item \textbf{Language Generation}: High-quality language generation that is coherent, relevant, and syntactically correct is essential.
\end{itemize}

\subsection{Evaluation Techniques}
\begin{itemize}
    \item \textbf{Accuracy Metrics}: Precision, recall, and F1-score are traditional metrics used to evaluate the correctness of answers.
    \item \textbf{Robustness Testing}: Assessing the model's performance under varying conditions and with adversarial inputs.
    \item \textbf{Fairness and Bias}: Evaluating the model for potential biases and ensuring that it provides equitable answers across different demographics.
    \item \textbf{Toxicity and Safety}: Ensuring that the model does not generate harmful or inappropriate content.
    \item \textbf{Human Feedback}: Incorporating human evaluations to assess the relevance and quality of the responses.
\end{itemize}

\citep{BERT2019, GPT22020, ChatGPT2022}.

\section{Challenges in Complex QA}
There are multiple challenges in the domain of complex QA which include but are not limited to:

\subsection{Domain Adaptation}
LLMs are typically trained on broad datasets, which can make them less effective for domain-specific questions. Techniques like fine-tuning and transfer learning are used to adapt models to specific domains.
\citep{TransferLearningSurvey2019}.

\subsection{Decomposition and Efficient Multi-step QA}
Complex questions often require decomposition into simpler sub-questions. The challenge is to manage and integrate the answers from these sub-questions into a coherent final answer. 
\citep{MultiHopQA2019}.

\subsection{Long Form and Non-Factoid QA}
Long-form QA involves generating detailed and structured answers. Non-factoid questions require explanations or opinions rather than factual answers, demanding sophisticated natural language understanding and generation.
\citep{NarrativeQA2018}.

\subsection{Safety and Multi-Sensitivity Data Protection}
Handling sensitive data and ensuring that the model does not breach privacy or produce harmful content is crucial.
\citep{DataPrivacy2021}.

\subsection{Multimodal Search}
Incorporating information from various modalities like text, images, and structured data to generate comprehensive answers.
\citep{MultimodalAI2020}.

\subsection{Hallucinations, Explainability, and Truthfulness}
LLMs can sometimes generate plausible but incorrect answers (hallucinations). Ensuring explainability and the truthfulness of answers is critical for user trust.
\citep{ExplainableAI2021}.

\section{Current Solutions and Research Trends}

\subsection{Hybrid LLM Architectural Patterns}
Integrating various architectures to combine strengths from different approaches remains a significant trend. For instance, combining LLMs with knowledge graphs or rule-based systems to improve accuracy and reliability.
\citep{HybridModels2020}.

\subsection{Training and Prompting Strategies}
Advanced techniques in training and prompting such as few-shot learning and active learning help improve model performance on specific tasks without extensive labeled data.
\citep{PromptingSurvey2021}.

\subsection{Active Human Reinforcement Learning Supervised with AI}
Combining machine learning models with real-time human feedback to iteratively improve model performance and address edge cases more effectively.
\citep{HumanInLoopAI2020}.

\subsection{Neuro-Symbolic and Structured Knowledge Grounding}
Integrating neural networks with symbolic reasoning to leverage the advantages of both paradigms, enhancing the model's ability to understand and generate accurate answers.
\citep{NeuroSymbolicAI2021}.

\subsection{Program Synthesis and Iterated Decomposition}
Using program synthesis techniques to parse and execute complex tasks iteratively can help break down questions into manageable components and generate accurate answers step-by-step.
\citep{ProgramSynthesis2019}.

\section{Conclusion}
The field of complex QA continues to evolve with innovations that seek to overcome the inherent limitations of existing models. The integration of hybrid architectures that leverage both LLMs and specialized techniques holds significant promise for addressing these challenges. As the AI community continues to advance the state-of-the-art through novel approaches and rigorous evaluations, the capability of models to handle complex, domain-specific questions is expected to improve significantly.

\bibliographystyle{plainnat}
\bibliography{prompt1_try1_4o}

\end{document}